\section{Research questions}

There are four key areas of focus that this research seeks to address. The first two are questions into the nature of biological process:

\begin{itemize}
\item  \textbf{Why did Y. pestis evolve a biphasic pathogenesis in humans?} 
Switching from pro-phagocytosis to an antiphagocytosis strategy seems to be crucial to the overall pathogenesis of Y. Pestis.  There is a simple argument based on the net proliferability both before and after neutrophils arrive on the scene of infection (see conclusion).  Also,  we could assume that Y. pestis replication is only slightly supercritical.  And it can be easily demonstrated that near critical branching process become drastically less likely to die out when the initial count is increased by even a few colony forming units.  If we see the initial within macrophage replication as setting up an initial population,  this effect offers further explanatory potential. 


\item \textbf{Is there a general advantage to "bursting" style transfer in successive intracellular infection methods?}
Natalie Komavora proposed that there could be a "flooding advantage" to bursting style pathogenesis.  That is,  bacteria or virus exiting their host infected cells all at once in large groups rather than one at a time or in dribs and drabs. Chapter 7 explores gives further mathematical support to this theory and discuses the conditions under which it might prove effective.
\end{itemize} 

\noindent
The second two are areas in which we tried to construct new models to describe particular dynamics:

\begin{itemize}
\item \textbf{A potential update to the standard BMVR model of viral replication to incorporate intracellular dynamics.}
Based on our analysis of the birth death catastrophe process, there is the possibility that one of the results could be encorporated into a rate of the BMVR.

\item \textbf{Efforts to describe variable rates of evolution with Markov chains}\par
The standard birth-death process has been used for decades to model the emergence of species over geological time periods. It may however be criticised for being consistent with the outdated doctrine of progressive gradualism. This states that evolutionary change is constant and happens gradually. In the 1970s,  S. J Gould proposed the counter theory of punctuated equilibria: evolution taking place rapidly over proportionally short periods which are intervened by longer times of relative morphological stasis.\par
There is also the controversial question of why large clades (containing many species) appear to generally have experienced an initial period of substantial diversification. Some attribute this observation to a statistical consequence of standard birth-death dynamics (push of the past, and pull of the present).\par
This work presents two models which seek to challenge this argument, and give mathematical basis to the punctuated equalibria theoretical paradigm. Notably, there is potential significance here in understanding the emergence of Y. Pestis.


\end{itemize}
